% =====================================================================
%  第 5 章  当前进度与后续计划
% =====================================================================

\section{当前进度与后续计划}
\label{sec:progress}

\subsection{已完成的工作}

\begin{enumerate}
  \item \textbf{选题论证}。确定了研究对象（极值图论的不变量不等式）、核心判断（瓶颈在筛选端）、以及三个子问题的划分。
  \item \textbf{智能体设计}。给出了闭环结构、四个阶段的分工、LLM 的四处职责、工具层接口与轨迹字段。这部分是本章之前的内容，属于设计而非实现。
  \item \textbf{基础环境}。搭好了 Python 环境、密钥管理（通过 \texttt{.env}，不写进代码）、Git 版本控制，以及本文档的 LaTeX 工程。
  \item \textbf{一个最小可运行的智能体原型}。在确定选题之前，先写了一个极简的 ReAct 循环（思考—行动—观察），用于理解智能体如何工作。它只有两个玩具工具（查天气、查景点），与本课题的数学内容无关，但验证了“LLM 能根据工具返回的结果决定下一步”这一基本机制。
\end{enumerate}

需要说明的是，第 4 项只是一个学习用的原型，不是本课题的智能体。本课题的系统尚未开始编码。

\subsection{尚未开始的部分}

\begin{center}
\small
\begin{tabular}{@{}llp{6.6cm}@{}}
\toprule
阶段 & 内容 & 产出 \\
\midrule
一 & 图样本与不变量计算 & 一张表：每行一张图，每列一个不变量 \\
二 & 猜想生成 & 给定数据，产出候选形式与系数 \\
三 & 反例搜索 & 对每条候选，在预算内搜索反例 \\
四 & 统计筛选 & 计算 $p$ 值、做多重检验校正 \\
五 & 端到端实验 & 完整的轨迹日志与对照实验数据 \\
六 & 分析与撰写 & 实验图表、报告后半部分 \\
\bottomrule
\end{tabular}
\end{center}

\subsection{时间安排}

\begin{center}
\small
\begin{tabular}{@{}llp{5.2cm}l@{}}
\toprule
日期 & 内容 & 产出 & 状态 \\
\midrule
11/19 & 选题论证、报告框架 & 本报告前两章 & 完成 \\
11/20 & 图样本与不变量 & 不变量数据表 & 待办 \\
11/21 & 猜想生成 & 首批候选 & 待办 \\
11/22 & 反例搜索、闭环跑通 & 第一份轨迹日志 & 待办 \\
11/23 & 统计筛选与对照实验 & 实验数据与图表 & 待办 \\
11/24 & 幻灯片与报告补充 & 汇报材料 & 待办 \\
11/25 & 中期汇报 & — & — \\
\bottomrule
\end{tabular}
\end{center}

\subsection{已知的风险}

\begin{center}
\footnotesize
\begin{tabular}{@{}p{3.6cm}p{10.4cm}@{}}
\toprule
风险 & 应对 \\
\midrule
不变量计算太慢 & 色数是 NP-难的，大图上算不准。方案：小图精确算，中等图用上下界代替并明确标注，大图只做采样。宁可结论保守，不假装精确。 \\
\addlinespace
LLM 调用不稳定 & 所有结果落盘缓存，重跑读缓存；汇报用录屏备份，不做现场演示。 \\
\addlinespace
时间不够 & 优先保证“闭环能跑通、有轨迹可展示”，统计筛选的对照实验可以缩减规模。 \\
\addlinespace
结果不理想 & 这本身是结论。如实报告失败模式，比编造成功更有价值。 \\
\bottomrule
\end{tabular}
\end{center}
